\documentclass[11pt,a4paper]{article}

\usepackage[utf8]{inputenc}

\usepackage[T1]{fontenc}

\usepackage[polish]{babel}

\usepackage{lmodern}

\usepackage{geometry}

\geometry{a4paper,margin=27mm}

\usepackage{microtype}

\usepackage[hidelinks]{hyperref}

% Portfolio: MSC-AR-P27; source: prospective ARC-XR research scope.

\title{Odporność adaptacyjnego HUD na błędy estymacji obciążenia poznawczego}

\author{Mateusz Pomianek}

\date{Wersja 0.3.0, 21 listopada 2026}

\begin{document}

\maketitle

\begin{abstract}

Przedmiotem opracowania jest Odporność adaptacyjnego HUD na błędy estymacji obciążenia poznawczego. Czy adaptacyjny HUD zachowuje stabilność przy opóźnieniach i utracie sygnału obciążenia? W kontrolowanych zadaniach HUD testowane będą różne poziomy opóźnienia, utraty danych i błędów estymacji obciążenia. Głównym punktem końcowym jest dodatkowe błędy zadań i liczba zbędnych zmian interfejsu względem warunku referencyjnego. Opis ma charakter wstępny i nie przedstawia niezweryfikowanych wyników.

\end{abstract}

\section{Wprowadzenie}

Rozwój adaptacyjnych systemów XR stwarza możliwość badania interakcji człowieka z informacją przestrzenną, ale samo zastosowanie AR, VR lub modeli sztucznej inteligencji nie przesądza o przewadze poznawczej czy ergonomicznej. W rozpatrywanym zagadnieniu, Odporność adaptacyjnego HUD na błędy estymacji obciążenia poznawczego, kluczowe jest rozdzielenie właściwości systemu od zmienności zachowania użytkownika. Porównanie należy zaprojektować tak, aby mierzyć jednoznaczne zjawisko, a nie tylko ogólną ocenę technologii.

\section{Cel pracy i pytanie badawcze}

Czy adaptacyjny HUD zachowuje stabilność przy opóźnieniach i utracie sygnału obciążenia? Projekt wymaga krytycznego przeglądu aktualnego stanu badań oraz zdefiniowania nowości w stosunku do innych tematów portfolio XR.

\section{Zakres i odrębny wkład naukowy}
Robustność dynamicznego HUD względem kontrolowanych błędów estymatora obciążenia i opóźnienia sygnału. Wpływ standardowej adaptacji wobec HUD statycznego to P01. Odporność przy konkretnych awariach wejścia, nie przyczynowy efekt wprowadzenia adaptacji w P01.
\section{Hipotezy badawcze}
\paragraph{H1.} Zabezpieczony filtr lub watchdog zmian HUD zmniejsza nadmiarowe błędy przy opóźnionym i brakującym sygnale względem prostego reaktywnego regulatora.
\paragraph{H2.} Watchdog redukuje liczbę zbędnych zmian HUD na minutę bez niewspółmiernego obniżenia dostępności wymaganej pomocy.
\paragraph{H3.} Korzyść jest obserwowana na całych niezależnych scenariuszach, a nie jedynie na wielokrotnie zakłóconej tej samej trajektorii.
\section{Projekt badania i plan eksperymentów}
\subsection{A. Operacjonalizacja i kwalifikacja danych}
Robustność dynamicznego HUD względem kontrolowanych błędów estymatora obciążenia i opóźnienia sygnału. Wpływ standardowej adaptacji wobec HUD statycznego to P01. Jednostka wnioskowania: scena/procedura z kontrolowaną trajektorią sygnału oraz odrębny sesyjny test u ludzi, jeśli zyska aprobatę. Przed badaniem należy określić źródła, wersje urządzeń, czas akwizycji, etykiety i niezależną referencję. Nie liczyć zależnych okien i klatek jako niezależnych prób.
\subsection{B. Eksperyment porównawczy}
Zamrozić adaptator i generator sygnału, a następnie wstrzykiwać opóźnienie, brak danych, skoki i dryf do jednakowych odtworzonych sesji. Wynik pierwszorzędowy: różnica dodatkowych błędów zadania względem warunku referencyjnego przy określonym profilu opóźnienia lub utraty sygnału. Zamrozić losowanie warunków, progi, mianownik punktu końcowego i kryteria wyłączeń bez dostępu do oceny końcowej.
\subsection{C. Kontrole i analiza odporności}
Ustalić referencję bez błędów, policyjny licznik zmian i rzeczywistą opóźnioną reakcję; nie przedstawiać syntetycznych zakłóceń jako pomiaru fizycznego. Z osobna raportować błędy i efekt stabilizacji, przy nieakceptowalnych warunkach system ma wracać do bezpiecznego statycznego HUD. Utrata sygnału, abstencja, przerwanie zadania i dane brakujące są odrębnymi kategoriami wyniku.
\subsection{D. Walidacja na niezależnym materiale}
Walidacja musi zachować podział po uczestniku, sesji, scenie lub zadaniu. Model porównań per niezależna scena z kontrastem odporny-regulator minus reaktywny; bootstrap po scenach i profilach zakłócenia. Raportować wyniki negatywne i rozbieżności z założeniami.
\section{Estymand, analiza statystyczna i decyzje}
różnica dodatkowych błędów zadania względem warunku referencyjnego przy określonym profilu opóźnienia lub utraty sygnału. Jednostka: scena/procedura z kontrolowaną trajektorią sygnału oraz odrębny sesyjny test u ludzi, jeśli zyska aprobatę. Model porównań per niezależna scena z kontrastem odporny-regulator minus reaktywny; bootstrap po scenach i profilach zakłócenia. H1 jest planowanym kontrastem głównym; H2 i H3 pozostają pomocnicze do zatwierdzenia hierarchii testów oraz kontroli wielokrotności. Wartości marginesów, mocy, minimalnego efektu istotnego praktycznie i liczebności mają wynikać z pilotażu, a nie z post hoc wyboru wyniku.
\section{Stan badań i granica nowości}
\begin{itemize}
\item Existing adaptive physiological feedback in AR means adapting an interface is not by itself a novel contribution. \cite{xrprior1}.
\item Physiology-driven VR adaptation has been demonstrated; independent behavioural gains require direct controlled trials. \cite{xrprior2}.
\end{itemize}
Odporność przy konkretnych awariach wejścia, nie przyczynowy efekt wprowadzenia adaptacji w P01. Odrębność na tle opublikowanych metod i pozostałych artykułów wymaga osobnego audytu.
\section{Warunki bezpieczeństwa, sprzętu i zgodności}
To prospektywna koncepcja, nie zgoda na eksperyment, prerejestracja lub wynik empiryczny. Praca z uczestnikami, interwencje stresowe i pomiar biometryczny wymagają przeglądu etycznego, świadomej zgody i ochrony danych. Quest 3 w wersji standardowej nie posiada eye trackingu, a Tobii Eye Tracker 5 jest urządzeniem ekranowym; nie zastępować jednego stanowiska drugim.
\section{Materiał i metody}

W kontrolowanych zadaniach HUD testowane będą różne poziomy opóźnienia, utraty danych i błędów estymacji obciążenia. Wymagane jest zdefiniowanie jednostki analizy, zasad doboru przypadków, wersji scenariusza i źródeł pomiaru. Rozkład warunków oraz reguły wykluczenia prób należy określić przed pomiarem. W badaniach z udziałem ludzi niezbędna jest uprzednia właściwa ocena etyczna i świadoma zgoda. W testach wyłącznie symulacyjnych należy oddzielić prawdę symulacyjną od obserwacji empirycznej.

\subsection{Wynik główny i kryteria jakości}

Dodatkowe błędy zadań i liczba zbędnych zmian interfejsu względem warunku referencyjnego. Pomiar powinien zachowywać jawne jednostki, znaczniki czasu i reguły oceny poprawności. Przy modelach predykcyjnych stosować rozłączne partycje oceny; przy eksperymentach interwencyjnych oddzielić skuteczność predykcji od przyczynowego wpływu interfejsu.

\section{Plan analizy}

Analiza funkcji degradacji względem parametrów błędu i reguły bramkowania jakości, w porównywalnych scenach. W wynikach należy raportować wielkość efektu lub błąd estymacji z przedziałem niepewności, liczbą obserwacji i odsetkiem odrzuconych prób. Analizy dodatkowe wymagają jasnego oznaczenia jako eksploracyjne. Nie wybierać progów decyzyjnych ani modelu na podstawie obserwacji ze zbioru testowego.

\section{Wyniki}

W obecnej wersji nie raportuje się wyników pomiarów, tabel statystycznych ani parametrów wyznaczonych z nieudokumentowanych danych. To kierunek badawczy, a nie zatwierdzony protokół badania ani zgłoszony artykuł.

\section{Dyskusja}

Interpretacja przyszłych wyników musi uwzględnić alternatywne mechanizmy i zależność wniosków od jakości danych, warunków zadania oraz ograniczeń sprzętu. Wartość potencjalnego rozwiązania zależy nie tylko od trafności modelu, lecz także od jego opóźnienia, stabilności i wpływu na zachowanie człowieka. Wyniki testów syntetycznych nie mogą być automatycznie uogólniane na populację uczestników.

\section{Ograniczenia badania}

Kontrolowany błąd syntetyczny pozwala testować odporność, ale nie potwierdza trafności czujników w badaniu na ludziach. Robustness adaptacji na błędnych sygnałach, zamiast zwykłej skuteczności statycznego vs adaptacyjnego HUD w MSC-AR-P01. Żadna część szkicu nie dowodzi przeprowadzenia zbierania danych, uzyskania zgód etycznych ani gotowości manuskryptu do zgłoszenia.

\section{Wnioski}

Przedstawiono problem, kontrast i sposób jego planowanej oceny. Rozstrzygnięcie postawionego pytania wymaga przygotowania protokołu, zweryfikowanych źródeł oraz poprawnie udokumentowanych wyników. Bez tych materiałów nie można sformułować konkluzji empirycznych.

\begin{thebibliography}{9}
\bibitem{xrprior1} Designing adaptive feedback systems for managing cognitive load in augmented reality (2025), Design Science 11, e49. \url{https://doi.org/10.1017/dsj.2025.10040}.
\bibitem{xrprior2} Towards Intelligent VR Training: A Physiological Adaptation Framework for Cognitive Load and Stress Detection (2025), UMAP 2025. \url{https://doi.org/10.1145/3699682.3727575}.
\end{thebibliography}
\end{document}

